\documentclass[12pt]{article}
\usepackage[utf8]{inputenc}
\usepackage[T1]{fontenc}
\usepackage[margin=1in]{geometry}
\usepackage{setspace}
\usepackage{xcolor}
\usepackage[colorlinks=true, urlcolor=blue!60!black, linkcolor=blue!60!black]{hyperref}
\setlength{\parindent}{0em}
\setlength{\parskip}{0.45em}
\renewcommand{\baselinestretch}{1.0}
\begin{document}
\begin{center}
{\Large \textit{Washington v.\ Davis}, 426 U.S.\ 229 (1976)\par}
\vspace{0.3em}
{\normalsize Case Brief --- One-Page Edition\par}
\end{center}

\textbf{Memory Jogger.} Disparate racial impact alone is not an equal protection violation; purpose is required (\href{https://supreme.justia.com/cases/federal/us/426/229/}{Justia}).

\textbf{Facts.} Black applicants challenged the D.C.\ police's ``Test 21,'' a verbal exam Black applicants failed at four times the white rate, without alleging intent. The Department actively recruited Black officers, and the test correlated with training results.

\textbf{Procedural History.} The D.C.\ Circuit, applying Title VII's \textit{Griggs} standard, held the test unconstitutional.

\textbf{Votes.} Reversed 7--2: White for the Court; Stevens concurred; Brennan, with Marshall, dissented on the statutory issue.

\textbf{Holding.} A neutral law is not unconstitutional solely because of disproportionate racial impact; discriminatory purpose must be shown, though impact is evidence.

\textbf{Reasoning.} Precedent required purpose, inferable from ``the totality of the relevant facts.'' Upgrading verbal skills is legitimate, and an impact rule would threaten countless neutral laws. Stevens: objective evidence often best proves intent.

\textbf{Examples.} \emph{Same-Side:} a veterans' preference favoring men (\textit{Feeney}). \emph{Opposite-Side:} a rezoning with procedural irregularities (\textit{Arlington Heights}). \emph{Fence-Sitter:} a bail algorithm using race-correlated proxies.

\textbf{Critique.} The intent rule ignores unconscious bias and rewards neutral-seeming discrimination; defenders say equal protection targets purposeful denial of equal status and leaves impact remedies to legislatures.

\textbf{Key Quotation.} ``Disproportionate impact is not irrelevant, but it is not the sole touchstone of an invidious racial discrimination forbidden by the Constitution.'' --- Justice White, 426 U.S.\ at 242.

\end{document}
